\documentclass[11pt,a4paper]{article}
\usepackage{amsmath,amssymb,amsthm,mathtools}
\usepackage{geometry}
\usepackage{hyperref}
\usepackage{booktabs}
\usepackage{enumitem}
\usepackage{xcolor}
\usepackage{newunicodechar}
\newunicodechar{ℝ}{\ensuremath{\mathbb{R}}}
\newunicodechar{∀}{\ensuremath{\forall}}
\newunicodechar{≥}{\ensuremath{\geq}}

\geometry{margin=1in}
\hypersetup{colorlinks=true,linkcolor=blue,citecolor=blue,urlcolor=blue}

\newtheorem{theorem}{Theorem}[section]
\newtheorem{lemma}[theorem]{Lemma}
\newtheorem{proposition}[theorem]{Proposition}
\newtheorem{definition}[theorem]{Definition}
\newtheorem{remark}[theorem]{Remark}

\title{\textbf{Metriplectic Completion of the Frohmanian Symplectic Tether:\\
Thermodynamic Structure and Lean Formalization for the Navier--Stokes Equations}}
\author{Benjamin Stanley Frohman\thanks{\textcopyright\ 2026 Benjamin Stanley Frohman. Distributed under the Apache License 2.0. Part of the Frohmanian Symplectic Tether program.}\\
Independent Researcher, Texas, United States\\
ORCID: \href{https://orcid.org/0009-0006-7068-3718}{0009-0006-7068-3718}\\
\href{mailto:frohmanbenjamin@gmail.com}{frohmanbenjamin@gmail.com}}
\date{Version 1.1, 10 November 2026}

\sloppy
\emergencystretch=3em
\begin{document}
\maketitle

\begin{abstract}
We construct a metriplectic completion of the Frohmanian Symplectic Tether on the coadjoint orbits of the volume-preserving diffeomorphism group \(\mathrm{SDiff}(\mathbb{T}^3)\). Building on Arnold's Lie--Poisson structure for ideal incompressible fluids and the modern metriplectic 4-bracket formalism of Morrison and collaborators for dissipative continuum systems, we introduce a dissipative metric bracket that restores full thermodynamic consistency (energy conservation together with nonnegative entropy production) while addressing the obstruction, already formalized in Lean~4, that a purely symmetric tether kernel yields a vanishing damping coefficient. In Lean~4 we formalize only the abstract algebraic structure of a metriplectic system (as a structure whose axioms are fields) and prove three algebraic identities about it: energy conservation \(\dot H=0\), entropy production \(\dot S=(S,S)\ge0\), and the symmetric-swap form of the damping obstruction. The checked module \texttt{MetriplecticNS\_checked.lean} compiles with no \texttt{sorry} and its theorems depend only on Lean's standard axioms (\texttt{propext}, \texttt{Classical.choice}, \texttt{Quot.sound}). The Navier--Stokes part of the Lean file is a placeholder (state type \texttt{Unit}, both brackets identically zero): Lean does \emph{not} prove, and this paper does not claim to prove, recovery of the Navier--Stokes or Navier--Stokes--Fourier equations; that recovery is cited background from the metriplectic literature. Energy degeneracy of the viscous bracket is an explicit assumption, and the simple viscous bracket written in Section~5 does not satisfy it. This is not a solution of the Clay Millennium problem; implications for conditional regularity are discussed only qualitatively.
\end{abstract}

\tableofcontents

\section{Introduction}
The Euler equations of ideal incompressible hydrodynamics possess a rich geometric structure: they are the geodesic equations on the infinite-dimensional Lie group of volume-preserving diffeomorphisms equipped with the \(L^2\) metric (Arnold, 1966). In the Hamiltonian formulation this is realized as Lie--Poisson dynamics on the dual of the Lie algebra of divergence-free vector fields. The addition of viscosity produces the Navier--Stokes equations, which are dissipative. A systematic way to incorporate dissipation while preserving the Hamiltonian skeleton and guaranteeing thermodynamic consistency (the first and second laws) is the metriplectic formalism.

A metriplectic system consists of a Poisson bivector (or Poisson bracket) generating the reversible dynamics and a symmetric positive-semidefinite metric tensor (or metric bracket) generating the irreversible dynamics. The modern 4-bracket formulation of Morrison and Updike unifies earlier binary-bracket approaches and provides an algorithmic construction of thermodynamically consistent models, including the Navier--Stokes--Fourier system and various multiphase extensions.

Independently, the Frohmanian Symplectic Tether has been introduced as a symmetric quadratic metric correction to the Arnold Lie--Poisson bracket on \(\mathrm{SDiff}(\mathbb{T}^3)\). While the tether satisfies several desirable geometric conditions (including controllable negative feedback at vorticity maxima), a kernel-checked theorem establishes a conditional obstruction: \emph{if} the kernel is symmetric and satisfies the degeneracy condition with respect to the kinetic-energy Hamiltonian (condition (C2), which is a hypothesis and is not proved for the tether), the damping coefficient vanishes identically. Consequently the pure geometric layer cannot by itself supply a positive dissipation rate.

The present work proposes adjoining an independent dissipative metric bracket to the tether, states the conditions such a bracket must satisfy, and formalizes in Lean~4 the abstract algebraic consequences of those conditions. The recovery of the classical viscous terms (or their Fourier-law counterparts) from a metriplectic bracket is known from the literature cited in Section~3; it is not reproved or formalized here.

\section{Arnold Lie--Poisson Structure for Ideal Fluids}
Let \(\mathfrak{g} = \mathfrak{X}_{\mathrm{div}}(\mathbb{T}^3)\) be the Lie algebra of smooth divergence-free vector fields on the three-torus. Its dual \(\mathfrak{g}^*\) may be identified with the space of vorticity fields via the \(L^2\) pairing. The Lie--Poisson bracket of two functionals \(F,G\) of the vorticity \(\omega\) is
\begin{equation}
\{F,G\}(\omega) = \int_{\mathbb{T}^3} \omega \cdot \Bigl(
\frac{\delta F}{\delta\omega}\times\frac{\delta G}{\delta\omega}\Bigr)\,dx.
\end{equation}
The kinetic-energy Hamiltonian
\begin{equation}
H(\omega) = \frac12\int_{\mathbb{T}^3}|u|^2\,dx,\qquad
u = \mathrm{curl}^{-1}\omega
\end{equation}
(the Biot--Savart operator) generates the Euler equations
\begin{equation}
\partial_t\omega + \mathcal{L}_u\omega = 0.
\end{equation}
The total helicity and the Casimir functionals associated with the coadjoint orbits are conserved.

\section{Metriplectic Structure of Dissipative Fluids}
A metriplectic system on a phase space with coordinates \(\xi\) is defined by a skew-symmetric Poisson bracket \(\{\,\cdot\,,\,\cdot\,\}\) and a symmetric metric bracket \((\,\cdot\,,\,\cdot\,)\) satisfying, for all functionals \(F,G\),
\begin{align}
\{F,G\} &= -\{G,F\},\qquad (F,G)=(G,F) && \text{(skew-symmetry, symmetry)},\\
\{F,S\} &= 0 && \text{(entropy is a Casimir)},\\
(F,H) &= 0 && \text{(energy degeneracy)},\\
(F,F) &\ge 0 && \text{(positive semidefiniteness)}.
\end{align}
In particular the entropy production is \((S,S)\ge 0\). (Version~1.0 wrote ``\((S,H)\ge0\)''; that was an error, since \((S,H)=(H,S)=0\) by energy degeneracy and symmetry.)
The evolution of an arbitrary observable is
\begin{equation}
\dot F = \{F,H\} + (F,S).
\end{equation}
In the 4-bracket formulation the metric bracket is recovered by contraction of a 4-linear form \((F,K;G,N)\) against the generators \(H\) and \(S\).

For the compressible Navier--Stokes--Fourier system the dynamical variables are mass density \(\rho\), momentum density \(m=\rho u\) and entropy density \(\sigma=\rho s\). The Hamiltonian is the total energy
\begin{equation}
H[\rho,m,\sigma] = \int\Bigl(\frac{|m|^2}{2\rho}+\rho\,U\bigl(\rho,\tfrac{\sigma}{\rho}\bigr)\Bigr)\,dx,
\end{equation}
the entropy is the Casimir \(S=\int\sigma\,dx\), and the Poisson bracket is the standard non-canonical fluid bracket. The metriplectic 4-bracket is constructed from the Kulkarni--Nomizu product of two positive-semidefinite operators that encode viscosity \(\mu\) and thermal conductivity \(\kappa\). In those works \cite{MorrisonUpdike2024,BarhamMorrisonZaidni2024} contraction recovers the viscous stress and Fourier heat flux, and the algebraic symmetries guarantee
\begin{equation}
\dot H=0,\qquad\dot S=\int\frac{\mu}{T}|\nabla u|^2+\frac{\kappa}{T^2}|\nabla T|^2\,dx\ge0.
\end{equation}

An analogous construction is available for the incompressible case, in which the metric bracket produces the term \(\nu\Delta\omega\) (or \(\nu\Delta u\)); the energy lost by the fluid is then carried by an internal-energy or entropy variable, so that the \emph{total} energy is conserved. All of this paragraph and the preceding one is cited background; none of it is formalized in this work.

\section{The Frohmanian Tether and the Damping Obstruction}
The Frohmanian tether kernel \(T\) is a symmetric bilinear form on functionals of the vorticity,
\begin{equation}
T(F,G)(\omega) = -\kappa\,\|\omega\|_{L^2}^2\int_{\mathbb{T}^3}
\bigl(\Pi_u\tfrac{\delta F}{\delta\omega}\bigr)\cdot
\bigl(\Pi_u\tfrac{\delta G}{\delta\omega}\bigr)\,dx,\qquad \kappa>0,
\end{equation}
where \(\|\omega\|_{L^2}^2=\int|\omega|^2\,dx\) is the global enstrophy (this is the form formalized in the companion Lean development; version~1.0 wrote a pointwise weight \(|\omega(x)|^2\) inside the integral, which is a different kernel) and \(\Pi_u\) is the \(L^2\)-orthogonal projector onto the complement of the velocity field recovered from \(\omega\). Because of the factor \(-\kappa<0\), the kernel is negative semidefinite: \(T(F,F)\le0\) for every \(F\). It therefore has the opposite sign to a metric bracket, which must satisfy \((F,F)\ge0\). The three tether conditions are:
\begin{enumerate}[label=(C\arabic*)]
\item coadjoint invariance (under suitable hypotheses; not proved for \(T\) in general),
\item degeneracy with respect to the kinetic-energy Hamiltonian, \(T(F,H)=0\) for all \(F\) (a hypothesis; \textbf{not proved} for \(T\)),
\item controllable negative feedback, \(T(F,F)=-\kappa\|\omega\|^2\int|\Pi_u\delta F/\delta\omega|^2\,dx\le0\).
\end{enumerate}
The Lean~4 formalization verifies the following conditional statement: for any symmetric kernel obeying (C2), that is \(T(M,H)=0\) for every \(M\), one also has
\begin{equation}
T(H,M)=0\qquad\text{for every functional }M.
\end{equation}
(Version~1.0 described \(T(M,H)=0\) as the verified conclusion; it is in fact the hypothesis (C2).) Consequently, \emph{if} (C2) holds, the pure tether supplies no net damping rate along the energy direction, and the geometric layer alone cannot close a regularity argument that requires a positive dissipative coefficient.

\section{Metriplectic Tether Completion}
We complete the structure by adjoining an independent dissipative metric bracket \(B_{\mathrm{visc}}\) that encodes the physical viscosity (and, when desired, thermal conductivity). The full metriplectic bracket is
\begin{equation}
(F,G)_{\mathrm{full}} = T(F,G) + B_{\mathrm{visc}}(F,G).
\end{equation}
For this to be a metriplectic bracket we need the following, which we state as \textbf{assumptions} (they are not proved here):
\begin{enumerate}[label=(A\arabic*)]
\item \(T(F,H)=0\) for all \(F\), i.e.\ (C2) for the tether;
\item \(B_{\mathrm{visc}}\) is symmetric and \(B_{\mathrm{visc}}(F,H)=0\) for all \(F\) (energy degeneracy of the viscous bracket);
\item \(B_{\mathrm{visc}}(F,F)\ge -T(F,F)\) for all \(F\), so that \((F,F)_{\mathrm{full}}\ge0\).
\end{enumerate}
Under (A1)--(A3), energy is conserved and the entropy production is
\begin{equation}
\dot S = (S,S)_{\mathrm{full}} = B_{\mathrm{visc}}(S,S) + T(S,S) \;\le\; B_{\mathrm{visc}}(S,S),
\end{equation}
which is \(\ge0\) by (A3). The tether can only lower the entropy production supplied by \(B_{\mathrm{visc}}\). (Version~1.0 wrote \(\dot S=B_{\mathrm{visc}}(S,H)\ge0\); by (A2) that quantity is \(0\), and the correct expression is the one above.)

A natural candidate in the incompressible setting is
\begin{equation}
B_{\mathrm{visc}}(F,G) = \nu\int_{\mathbb{T}^3}
\nabla\Bigl(\tfrac{\delta F}{\delta\omega}\Bigr):\nabla\Bigl(\tfrac{\delta G}{\delta\omega}\Bigr)\,dx .
\end{equation}
It is symmetric and positive semidefinite, but it does \emph{not} satisfy (A2) on functionals of \(\omega\) alone. Indeed \(\delta H/\delta\omega=\psi\), the divergence-free, mean-zero stream function with \(-\Delta\psi=\omega\), and
\[
B_{\mathrm{visc}}(H,H)=\nu\int_{\mathbb{T}^3}|\nabla\psi|^2\,dx>0\qquad\text{whenever }\omega\neq0 .
\]
So with this candidate the fluid's kinetic energy is not conserved, which is the physically correct behaviour of viscous flow, but it means the candidate is not an energy-degenerate metric bracket. A bracket satisfying (A2) requires an additional internal-energy or entropy variable that absorbs the dissipated energy, as in the constructions cited in Section~3. That construction is not carried out here. Whether (A1) holds is open.

If (A1)--(A3) hold, the resulting system is thermodynamically consistent and retains the geometric tether for possible use in conditional regularity arguments. That such a system reproduces the Navier--Stokes equations is not shown here.

\section{Thermodynamic Consistency Theorems}
\begin{theorem}[Energy conservation]
\label{thm:energy}
Let \(\dot F = \{F,H\} + (F,S)\). Then
\[
\dot H = \{H,H\} + (H,S) = 0,
\]
where \(\{H,H\}=0\) by skew-symmetry of the Poisson bracket and \((H,S)=0\) by the null-vector condition on the metric bracket (together with symmetry).
\end{theorem}

\begin{theorem}[Entropy production]
\label{thm:entropy}
\[
\dot S = \{S,H\} + (S,S) = (S,S) \ge 0,
\]
since \(\{S,H\}=-\{H,S\}=0\) by skew-symmetry and the Casimir property of \(S\) (applied with \(F=H\)), and the metric bracket is positive semi-definite.
\end{theorem}

\begin{theorem}[Damping obstruction for symmetric kernels]
\label{thm:obstruction}
Any symmetric bilinear form \(B\) that is degenerate on the Hamiltonian (\(B(F,H)=0\) for all \(F\)) necessarily satisfies \(B(H,F)=0\) for all \(F\). Consequently, if a purely symmetric tether kernel obeys condition (C2), it supplies a vanishing damping coefficient along the Euler vector field.
\end{theorem}

\begin{remark}[Recovery of Navier--Stokes--Fourier: cited, not proved here]
\label{thm:recovery}
According to the cited literature, under the standard choices of Hamiltonian, entropy, non-canonical Poisson bracket and the metriplectic 4-bracket constructed from the viscosity and conductivity operators (Morrison--Updike, Barham--Morrison--Zaidni), the evolution law recovers the compressible Navier--Stokes--Fourier equations in weak form, including the precise entropy-production identity
\[
\dot S = \int\frac{\mu}{T}|\nabla u|^2 + \frac{\kappa}{T^2}|\nabla T|^2\,dx \ge 0.
\]
We do not prove this statement, and it is not formalized in Lean.
\end{remark}

The proofs of Theorems~\ref{thm:energy}--\ref{thm:obstruction} are purely algebraic and are formalized in the accompanying Lean module (\texttt{MetriplecticNS\_checked.lean}). Remark~\ref{thm:recovery} is background from the literature, not a result of this paper.

\section{Lean~4 Formalization}
The abstract algebraic structure is encoded as a Lean structure (with \texttt{Functional State := State → ℝ})
\begin{verbatim}
structure MetriplecticSystem (State : Type u) where
  poisson : Functional State → Functional State → Functional State
  metric  : Functional State → Functional State → Functional State
  H : Functional State
  S : Functional State
  poisson_skew : ∀ F G, poisson F G = fun s => - (poisson G F s)
  metric_sym   : ∀ F G, metric F G = metric G F
  metric_pos   : ∀ F s, metric F F s ≥ 0
  casimir_S    : ∀ F s, poisson F S s = 0
  null_H       : ∀ F s, metric F H s = 0
\end{verbatim}
\noindent\textbf{What Lean proves.} The checked module \texttt{MetriplecticNS\_checked.lean} proves \texttt{energy\_conserved} (\(\dot H=0\)), \texttt{entropy\_production} (\(\dot S\ge0\)) and \texttt{tether\_damping\_obstruction} (Theorem~\ref{thm:obstruction}). It compiles with Lean \texttt{v4.30.0-rc1} and Mathlib \texttt{v4.30.0-rc1} with no errors and no \texttt{sorry}, and \texttt{\#print axioms} reports only \texttt{propext}, \texttt{Classical.choice} and \texttt{Quot.sound} for all three theorems. The original module \texttt{MetriplecticNS.lean} is included unchanged for provenance; it does not compile (the entropy proof cited the Casimir condition at \(F=S\), giving \(\{S,S\}=0\) instead of the needed \(\{S,H\}=0\), and two doc comments were misplaced). The checked module changes only those proof steps and comments, not the statements.

\noindent\textbf{What Lean does not prove.} The Navier--Stokes section of the Lean file is a placeholder: the state type is \texttt{Unit} and both brackets are identically zero. No integral functionals, functional derivatives, Sobolev spaces, integration by parts, viscous term, 4-bracket or recovery of the Navier--Stokes or Navier--Stokes--Fourier equations is formalized. (Version~1.0 said the specialization to fluid fields and the recovery of the viscous term by integration by parts were formalized; that was not true.) Likewise, inserting the tether kernel as a summand of the metric bracket appears in the Lean file only as a comment. The obstruction theorem by itself says nothing about \(T(S,S)\); the statement about entropy production in a split bracket is the inequality of Section~5, under assumptions (A1)--(A3).

\section{Conclusion and Open Problems}
The metriplectic completion supplies the thermodynamically consistent dissipative structure that was missing from the pure geometric tether. Because of the obstruction already proved for any symmetric kernel satisfying the kinetic-energy degeneracy, the tether itself cannot furnish the analytic dissipation rate required for a regularity proof. The rate must come from the classical viscous bracket (or from a carefully constructed higher-order geometric correction that evades the obstruction). 

Open questions include:
\begin{itemize}
\item construction of a non-symmetric or higher-order geometric correction that yields a positive rate while preserving coadjoint equivariance;
\item proof or refutation of (C2) for the tether kernel, and construction of an energy-degenerate viscous bracket satisfying (A2) and (A3);
\item extension of the Lean formalization from the abstract setting to functionals on three-dimensional Sobolev spaces, including recovery of the Navier--Stokes equations;
\item quantitative comparison of the metriplectic discretizations (structure-preserving schemes) with the conditional majorant techniques already developed for the tether.
\end{itemize}
The present work provides the algebraic and formal scaffolding on which those investigations can proceed. Nothing here proves global regularity of the three-dimensional Navier--Stokes equations, and this is not a solution of the Clay Millennium problem.

\section*{Changes from version 1.0}
\begin{itemize}
\item Abstract, Introduction and Section~7: removed the claim that Lean proves recovery of the continuum equations; stated exactly what \texttt{MetriplecticNS\_checked.lean} proves and that the Navier--Stokes part of the Lean file is a placeholder.
\item Section~3: the entropy-production condition \((S,H)\ge0\) is corrected to \((S,S)\ge0\) (from \((F,F)\ge0\)); skew-symmetry and symmetry are stated.
\item Section~4: the tether formula uses the global enstrophy factor as formalized; the sign \(-\kappa\) makes \(T\) negative semidefinite, which is stated; (C2) is marked as a hypothesis that is not proved, and the Lean result is stated correctly as \(T(H,M)=0\) assuming (C2).
\item Section~5: \(\dot S=B_{\mathrm{visc}}(S,H)\) corrected to \(\dot S=(S,S)_{\mathrm{full}}\); energy degeneracy of \(B_{\mathrm{visc}}\) stated as an assumption, with a proof that the displayed candidate does not satisfy it; the ``recovers Navier--Stokes'' claim removed.
\item Section~6: the recovery theorem is now a remark citing the literature.
\item The Lean status, toolchain and axiom output are recorded; copyright, ORCID, affiliation and program note added.
\end{itemize}

\begin{thebibliography}{9}
\bibitem{Arnold1966}
V.~I. Arnold.
\newblock Sur la géométrie différentielle des groupes de Lie de dimension infinie et ses applications à l'hydrodynamique des fluides parfaits.
\newblock \emph{Ann. Inst. Fourier} \textbf{16} (1966), 319--361.

\bibitem{Morrison1986}
P.~J. Morrison.
\newblock A paradigm for joined Hamiltonian and dissipative systems.
\newblock \emph{Physica D} \textbf{18} (1986), 410--419.

\bibitem{Materassi2014}
M. Materassi.
\newblock Metriplectic algebra for dissipative fluids in Lagrangian formulation.
\newblock arXiv:1409.6542, 2014.

\bibitem{MorrisonUpdike2024}
P.~J. Morrison and M.~H. Updike.
\newblock Inclusive curvature-like framework for describing dissipation: metriplectic 4-bracket dynamics.
\newblock \emph{Phys. Rev. E} \textbf{109} (2024), 045202.

\bibitem{ZaidniMorrisonBenjelloun2024}
A. Zaidni, P.~J. Morrison and S. Benjelloun.
\newblock Thermodynamically consistent Cahn--Hilliard--Navier--Stokes equations using the metriplectic dynamics formalism.
\newblock arXiv:2402.11116, 2024.

\bibitem{BarhamMorrisonZaidni2024}
W. Barham, P.~J. Morrison and A. Zaidni.
\newblock A thermodynamically consistent discretization of 1D thermal-fluid models using their metriplectic 4-bracket structure.
\newblock arXiv:2410.11045, 2024.

\bibitem{Frohman2026}
B.~S. Frohman.
\newblock The Frohmanian Tether on \(\mathrm{SDiff}(\mathbb{T}^3)\): a metriplectic correction, its canonicity, and an obstruction to damping.
\newblock Zenodo, 2026. DOI: 10.5281/zenodo.24120781.
\end{thebibliography}

\end{document}
